\chapter{Ball Flight Models and Trajectory Estimation}
\label{ch:flight}

A launch monitor connects three different problems: observing motion, estimating a launch state, and predicting the flight that follows.
An indoor carry value is principally a prediction from a short observed segment; an outdoor system may also observe much of the trajectory.
The same aerodynamic model can be used forward to predict carry or inversely to infer launch conditions from later observations.
Neither direction makes the model exact, and using one model in both directions does not provide independent validation.

\section{Equations of Motion and Coefficient Conventions}

Let $\vect{p}$ and $\vect{v}$ denote ball-centre position and velocity in a ground-fixed frame, and let $\vect{w}(\vect{p},t)$ be the local wind velocity.
Define air-relative velocity $\vect{v}_a=\vect{v}-\vect{w}$, airspeed $V=\lVert\vect{v}_a\rVert$, cross-sectional area $A=\pi R^2$, and dynamic pressure $q=\rho V^2/2$.
Radius $R$ and mass $m$ should describe the ball being modeled; values near 21.34\,mm and 45.93\,g are illustrative inputs, not identities for every ball.
A quasi-steady drag-and-Magnus model is
\begin{equation}
\label{eq:eom}
\dot{\vect{p}}=\vect{v},\qquad
m\dot{\vect{v}}=-qAC_D\frac{\vect{v}_a}{V}
                   +qAC_L\vect{e}_L+m\vect{g},\qquad
\vect{e}_L=\frac{\vect{\omega}\times\vect{v}_a}
                    {\lVert\vect{\omega}\times\vect{v}_a\rVert}.
\end{equation}
Here $\vect{\omega}$ is the physical spin vector and $C_L$ is a \emph{signed} coefficient in the displayed direction.
This convention assigns lift magnitude $qA|C_L|$; a formula using $\hat{\vect{\omega}}\times\hat{\vect{v}}_a$ instead has an additional factor equal to the sine of the angle between spin and airflow.
Coefficients fitted under one convention cannot be moved to the other without that conversion.
With the book's right-handed axes ($x$ forward, $y$ up, $z$ right), positive $z$ backspin and forward air-relative motion give upward lift when $C_L>0$.
The force law is an approximation to the mean aerodynamic force; it does not represent every unsteady or orientation-dependent force on a real dimpled ball.

The unit vectors in \cref{eq:eom} require explicit limiting rules.
At zero airspeed, use the model's zero-force limit rather than divide by $V$.
At zero spin or parallel spin and airflow, the cross-product direction is undefined.
A symmetric Magnus model sets that lift contribution to zero and requires its coefficient to approach zero consistently; an empirical fit measured away from these limits does not establish its own continuation to them.
For drag force $\vect{F}_D=-qAC_D\vect{v}_a/V$ with $C_D\geq0$, the air-relative power is nonpositive: $\vect{F}_D\cdot\vect{v}_a=-qAC_DV$.
It can nevertheless do positive work in the ground frame when a sufficiently fast tailwind accelerates the ball.
Lift in this model is perpendicular to air-relative velocity, not necessarily to ground-relative velocity.

Useful dimensionless inputs are Reynolds number $\mathrm{Re}=2RV/\nu$ and spin ratio $S=R\lVert\vect{\omega}\rVert/V$, with kinematic viscosity $\nu$.
The angle $\alpha$ between spin and airflow, ball construction and surface condition may also matter.
A model interface should therefore declare its dependencies and calibration domain, schematically
\begin{equation}
\label{eq:smits}
C_D=f_D(\mathrm{Re},S,\alpha;\vect{b}),\qquad
C_L=f_L(\mathrm{Re},S,\alpha;\vect{b}),
\end{equation}
where $\vect{b}$ denotes a versioned set of ball-specific parameters.
This is an interface definition, not a numerical fit attributed to a particular author.
Fits obtained with spin perpendicular to the airflow need an additional justified assumption before use at general $\alpha$.

Spin evolution needs its own equation, for example $I\dot{\vect{\omega}}=\vect{\tau}_{a}$ for an idealized ball with constant scalar moment of inertia $I$ and aerodynamic torque $\vect{\tau}_{a}$.
An empirical decay law $\dot{\vect{\omega}}=-k\vect{\omega}$ further assumes no change in axis direction from torque, and $k$ may depend on speed and spin.
Even when physical spin decays, spin ratio can rise as the ball slows.
For fixed $R$, nonzero spin and positive $V$,
\begin{equation}
\label{eq:spinratioevolution}
\frac{\dot S}{S}
=\frac{\vect{\omega}\cdot\dot{\vect{\omega}}}{\lVert\vect{\omega}\rVert^2}
-\frac{\dot V}{V}.
\end{equation}
Thus integrating $S$ with a physical-spin decay constant is generally inconsistent unless the airspeed term is included.
For a varying wind field, $\dot V=\vect{v}_a\cdot\dot{\vect{v}}_a/V$ uses $\dot{\vect{v}}_a=\dot{\vect{v}}-\partial_t\vect{w}-(\vect{v}\cdot\nabla)\vect{w}$, including the wind variation along the ball's path.

\section{What Aerodynamic Data Establish}

Bearman and Harvey's wind-tunnel study compared particular spinning, dimpled ball models~\cite{bearmanharvey}.
It motivates accounting for Reynolds number, spin and surface geometry; its tested dimple comparison is not a universal ranking of modern balls.
Smits and Smith~\cite{smitssmith} and Quintavalla~\cite{quintavalla} provide further modeling approaches.
Quintavalla's author description identifies six-term lift and drag models fitted to Indoor Test Range measurements over driver trajectories.
That is a calibration strategy, not evidence that any six-term polynomial is accurate outside its fitted conditions or intrinsically superior to wind-tunnel data.
This review does not prescribe numerical coefficients from an unverified transcription of those models; an implementation must establish the original equations and their domain before adopting them.

An implementer needs the actual coefficient data, normalization, units, fitted ball, validity range and spin-torque law, not just a model name.
Preserve source tables or reproducible digitization and report extrapolation beyond the calibration domain.
Fit aerodynamic parameters using independent launch observations and multiple trajectories with varied speeds and spins.
Otherwise a biased launch angle, unknown wind or a spin error can be absorbed into apparently plausible lift and drag coefficients.
US6186002 describes determination of coefficients from measured trajectories~\cite{us6186002}; a patent's disclosure supplies a method description, not independent evidence of accuracy or a legal clearance determination.

\begin{implication}
Version the coefficient data, spin law, atmosphere, wind, ground definition and numerical solver together.
Validate predicted trajectories on shots excluded from fitting.
Agreement with another simulator is a comparison of predictions; it does not establish that either prediction matches the physical flight.
\end{implication}

\section{Numerical Integration and the Landing Event}

Fourth-order Runge--Kutta is one possible integrator, but its order alone does not establish adequate carry accuracy.
Check convergence as the step is reduced, resolve coefficient transitions, and locate landing between integration steps.
Define landing as a descending intersection of the ball-centre trajectory with a declared height surface, excluding the launch-time root.
Returning to launch elevation is only one convention; tee height, uneven ground and a surface referenced to the ball's contact point require explicit treatment.
Bounce and roll require additional contact and terrain models and are not part of airborne carry.

Landing-time uncertainty matters even if the trajectory integrator is exact.
Let $a$ be a scalar uncertain input and let the landing event satisfy $G(\vect{p}(t_*,a),t_*,a)=0$.
For fixed-time sensitivity $Y=\partial\vect{p}/\partial a$, implicit differentiation gives
\begin{equation}
\label{eq:landingderivative}
\frac{dt_*}{da}=-\frac{G_{\vect{p}}Y+G_a}{G_{\vect{p}}\vect{v}+G_t},\qquad
\frac{d\vect{p}_*}{da}=Y+\vect{v}\frac{dt_*}{da}.
\end{equation}
All terms are evaluated at the same landing event.
The formula assumes differentiability and a transverse crossing, $G_{\vect{p}}\vect{v}+G_t\ne0$; a nearly tangential crossing is poorly conditioned.
For a vector of uncertain inputs, differentiate with respect to each input component in turn.
For an output $c(\vect{p}_*,t_*,a)$, also include the explicit terms $c_t\,dt_*/da+c_a$ along with $c_{\vect{p}}\,d\vect{p}_*/da$.
For downrange position $c=x_*$ with a fixed origin, those explicit terms vanish; the landing-time contribution is already included in $d\vect{p}_*/da$.

A vacuum trajectory makes the missing term easy to see.
For level launch and landing, $x=v_x t$, $y=v_y t-g_0t^2/2$, and $v_y>0$, the descending root is $t_*=2v_y/g_0$.
Downrange carry is $c=2v_xv_y/g_0$, so $\partial c/\partial v_y=2v_x/g_0$.
At fixed time, however, $\partial x/\partial v_y=0$.
Ignoring the change in landing time would therefore erase the entire vertical-launch-speed contribution to carry uncertainty in this simple example.

\section{Trajectory Estimation and Filtering}
\label{sec:ekf}

A useful physical state is $\vect{x}=(\vect{p},\vect{v},\vect{\omega})$, augmented only with additional parameters that the data can constrain.
For a fixed radar at position $\vect{s}$, range is $r=\lVert\vect{p}-\vect{s}\rVert$ and radial velocity is $\dot r=\vect{n}\cdot\vect{v}$, with $\vect{n}=(\vect{p}-\vect{s})/r$.
Angles require the radar's own axes; camera pixels require calibrated projection, pose and timing.
Different sensors therefore need different measurement functions, Jacobians and noise models, even when they share one state estimator.
Observation windows, processing delays, association of returns to ball or club, and correlated errors belong in that model; treating every output as a simultaneous independent point measurement overstates information.

Extended or unscented Kalman filtering and subsequent smoothing provide approximate state estimates conditional on the process and observation models~\cite{sarkkasvensson2023}.
Backward estimation to ball separation can use the observed arc, but its uncertainty grows with uncertain timing, model error and weak geometric information.
Including a spin-axis state does not make that axis observable.
A short arc can be consistent with multiple combinations of launch velocity, wind, lift coefficients and spin; a tight prior may select one combination without the observations independently identifying it.
Check sensitivity rank and conditioning, residual structure and performance on withheld data, rather than equating a small filter covariance with demonstrated accuracy.

For the K-LD7 example, use the speed-dependent measurement cadence documented in \cref{app:impl}; a roughly 29\,ms cycle at the stated 100\,km/h example does not supply 10--20 independent detections in 50\,ms.
Processing a complete burst can help only to the extent that it contains useful, correctly associated observations.
Neither a smoother nor a confidence weight repairs an incorrect timestamp or a missing observation direction.

\section{Indoor Extrapolation and Coupled Uncertainty}

Let $\vect{a}$ collect launch state, aerodynamic parameters, atmospheric inputs and ground geometry, and let $c(\vect{a})$ be a carry or lateral-distance prediction evaluated at landing.
For locally smooth behavior and sufficiently small uncertainty,
\begin{equation}
\label{eq:flightcovariance}
\operatorname{Var}(c)\approx J_c\Sigma_aJ_c^{\mathsf T},\qquad
J_c=\frac{dc}{d\vect{a}}.
\end{equation}
The covariance $\Sigma_a$ must retain correlations, including those introduced when launch and aerodynamic parameters are fitted together.
This expression covers uncertainty in the declared inputs, not every discrepancy between the chosen force law and reality.
Model discrepancy needs separate evidence from held-out trajectories; combine it probabilistically only after stating its dependence on the input errors.
For nonlinear, multimodal or threshold-sensitive predictions, propagate representative joint samples and report the resulting distribution rather than only a local Gaussian approximation.

There is consequently no universal conversion from 300\,rpm of spin error, two degrees of axis error, or a 0.01 change in drag coefficient to a fixed number of yards.
Report the baseline launch, ball model, atmosphere, perturbation and landing convention with any numerical sensitivity.
Likewise, tilting spin changes lift direction only under stated geometric and coefficient assumptions; lateral displacement integrates that changing force over the flight and is not a universal percentage of carry per degree.
An optical spin observation can constrain an indoor extrapolation, but its calibration and rotation ambiguity still need validation.
The evidence question is which uncertain inputs control this shot's prediction and how well they were observed, not which sensor family wins by definition.
